\textbf{\Large{Instrucciones:}}
\begin{enumerate}
    \item Este ensayo contiene \NumPreguntas{} preguntas de selección múltiple, cada una con \NumAlternativas{} opciones (\LetrasAlternativas), de las cuales solo una es correcta. Todas las preguntas tienen el mismo valor.
    \item Dispones de \Duracion{} (\Minutos{} minutos) para responder las \NumPreguntas{} preguntas.
    \item Las respuestas válidas solo se consideran en la hoja de respuestas. Marca una sola opción por pregunta.
    \item No se descuenta puntaje por respuestas erradas.
    \item Puedes usar este cuadernillo como borrador, pero solo se corregirá lo marcado en la hoja de respuestas.
    \item No se permite el uso de calculadora, teléfono celular ni de ningún otro dispositivo electrónico.
    \item Al terminar, entrega la hoja de respuestas y este cuadernillo.
\end{enumerate}

\ifconsideracionesmatematicas
\vspace{0.3cm}
\textbf{\large{En este ensayo se considerará lo siguiente:}}
\begin{enumerate}
    \item Las figuras que aparecen son solo indicativas.
    \item Los gráficos que se presentan están dibujados en un sistema de ejes perpendiculares, cuyo origen es el punto $(0,\ 0)$.
    \item El intervalo $[p,\ q]$ es el conjunto de todos los números reales mayores o iguales que $p$ y menores o iguales que $q$; el intervalo $]p,\ q]$ es el conjunto de todos los números reales mayores que $p$ y menores o iguales que $q$; el intervalo $[p,\ q[$ es el conjunto de todos los números reales mayores o iguales que $p$ y menores que $q$; y el intervalo $]p,\ q[$ es el conjunto de todos los números reales mayores que $p$ y menores que $q$.
    \item $\vec{v}=(a,\ b)$ es un vector que tiene su punto de inicio en el origen del plano cartesiano y su punto final en el punto $(a,\ b)$, a menos que se indique lo contrario.
    \item Se entenderá por dado común a aquel que posee 6 caras numeradas del 1 al 6 y, en el experimento de lanzarlo, sus caras son equiprobables de salir.
    \item En el experimento de lanzar una moneda, sus dos opciones son equiprobables de salir, a menos que se indique lo contrario.
\end{enumerate}

\vspace{0.2cm}
\textbf{\large{Símbolos matemáticos:}}\par\vspace{0.2cm}
\begin{tabular}{cl}
$<$ & es menor que \\ $>$ & es mayor que \\ $\leq$ & es menor o igual que \\ $\geq$ & es mayor o igual que \\
\tikz[baseline=0pt, line width=0.6pt]{\draw (0.3,0) -- (0,0) -- (0,0.28); \draw (0,0.11) -- (0.11,0.11) -- (0.11,0);} & ángulo recto \\
$\sim$ & es semejante con \\ $\neq$ & es distinto de \\ $\overline{AB}$ & trazo $AB$ \\ $AB$ & medida del trazo $AB$ \\
\end{tabular}
\fi
